% Part: proof-theory
% Chapter: cut-elimination
% Section: introduction

\documentclass[../../../include/open-logic-section]{subfiles}

\begin{document}

\olfileid{pt}{cut}{int}

\olsection{Introduction}

The \Cut{} rule is the rule
\[
\Axiom$\Gamma \fCenter \Delta, !A$
\Axiom$!A, \Pi \fCenter \Lambda$
\RightLabel{\Cut}
\BinaryInf$\Gamma, \Pi \fCenter \Delta, \Lambda$
\DisplayProof
\]
(The !!{formula}~$!A$ is called the \emph{cut formula}.)

\Cut{} was included in Gentzen's original sequent calculus~\Log{LK}.
Gentzen's \emph{Hauptsatz} (``main theorem'') states that every
sequent !!{provable} in~\Log{LK} is also !!{provable} without using
the \Cut{} rule, i.e., that \Cut{} can be ``eliminated'' from
\Log{LK}-proofs. Our systems \Log{G1c} and \Log{G3c} do not contain
\Cut, but the question can be raised for them as well: can \Cut{} be
eliminated from !!{proof}s that use \Cut{} plus the other rules of
\Log{G1c} (or those of \Log{G3c})? In our established terminology it
can be asked equivalently as: Is \Cut{} an admissible rule of
\Log{G1c} (or \Log{G3c}).

The cut elimination theorem is perhaps the most important and
celebrated result of proof theory. It first of all establishes that
\Log{G1c} and \Log{G3c} don't need a rule which you might, at first
blush, think is clearly required. For when you try to formalize actual
proofs, you'll notice that you need a way to deal with the use of
lemmas. Mathematicians like to prove results, and then appeal to these
results in proofs of other results. So you might first formalize a
proof of result $L$ (the lemma) as the !!{proof} of a sequent
$\Sequent L$. Then you formalize the proof of the second result~$R$,
the proof of which uses~$L$, as the !!{proof} of a sequent $L \Sequent
R$. How do you know that there is a proof of~$R$, i.e., !!a{proof} of
just $\Sequent R$? You need a way to combine the two !!{proof}s, and
the \Cut{} rule would allow you to do that.

We mentioned that \Log{G1c} and \Log{G3c} are complete, and so they
prove any sequent you might expect to be able to prove---all valid
sequents. This can be proved directly. But at the time Gentzen was
writing, only axiomatic derivation systems were known to be complete.
To show \Log{LK} complete, Gentzen provided a translation of
!!{derivation}s in the axiomatic system into \Log{LK}-!!{proof}s. This
translation used \Cut{} essentially. Cut elimination is not just an
important result for classical logic: many other logics have sequent
calculi. For some logics, a direct proof of completeness for their
sequent calculi is hard or impossible, so translations from other
systems (using \Cut) are the only way to show completeness. Then
proving cut elimination proof-theoretically is necessary to establish
that \Cut{} is superfluous and that the system without \Cut{} is
complete.

Cut elimination is also important because it can be applied to obtain
other, independently interesting results. Two cases we'll consider are
the Craig interpolation lemma and Herbrand' theorem.

The ideas used to prove cut elimination typically involve showing how
!!a{proof} that uses \Cut{} can be transformed into one without. These
transformations are usually ``local,'' i.e., we take a part of !!
a{proof} (usually a sub-!!{proof} ending in a \Cut{} inference) and
replace it with another sub-!!{proof} of the same sequent but in which
the \Cut{} inference has been removed or replaced by others. Such
arguments provide step-by-step algorithms for transforming !!{proof}s
with \Cut{} into ones without. This provides more information when cut
elimination procedures are used in applications. For instance, there
are model theoretic proofs of the interpolation lemma the form ``if
$!A \lif !B$'' is valid then an interpolant exists (never mind what
and interpolant is for the moment). A proof-theoretic argument using
cut elimination provides an algorithm that takes !!a{proof} of $!A
\lif !B$ and returns an interpolant. In contrast to the model
theoretic argument, the proof theoretic argument is
\emph{constructive}.

There are two popular ways to prove the cut elimination theorem. One
(going back to Gentzen) shows that you can remove top-most cuts in a
proof, and then just removes top-most cuts until none are left. The
other (originally due to Sch\"utte and Tait) focuses on the most
complex cuts in !!a{proof}, and successively removes cuts that are
maximal in the sense that no other cut has a cut !!{formula} of higher
depth. They each have their advantages and drawbacks. For some
systems, one approach works while another is hard to implement or even
impossible. So, we describe both approaches. We'll prove cut
elimination for~\Log{G3c}. In fact, we'll show not that \Cut{} can be
eliminated, but \emph{context-sharing cut}~\CutCS:
\[
\Axiom$\Gamma \fCenter \Delta, !A$
\Axiom$!A, \Gamma \fCenter \Delta$
\RightLabel{\Cut}
\BinaryInf$\Gamma \fCenter \Delta$
\DisplayProof
\]
If a sequent calculus allows weakening and contraction (either as
actual rules like in \Log{G1c} or as admissible rules like in
\Log{G3c}), \Cut{} can simulate \CutCS, and vice versa. We choose
\CutCS{} over \Cut because the first cut elimination strategy is
easier to describe for~\CutCS, and the second strategy only works
for~\CutCS.

\end{document}
